\documentclass{leveldoc}
\DocUses{foundations,slingshot,first_coil,first_magnet}
\DocFingerprint{60121b6799d04e86}

\begin{document}

\section{What the level asks}\label{sec:asks}

The game's text:
\leveldescription

Three positive ions of masses $\val{mOne}$, $\val{mTwo}$ and $\val{mFour}$ leave the source at
$(\val{xZero}, \val{yZero})$ together, along $+x$, with the same kinetic energy
$T_0 = \val{Tzero}$. Each must enter its own detector at the right of the board, the lightest
the highest: mass 1 the box at $\val{yDetLoOne} \le y \le \val{yDetHiOne}$, mass 2 at
$\val{yDetLoTwo} \le y \le \val{yDetHiTwo}$, mass 4 at $\val{yDetLoFour} \le y \le \val{yDetHiFour}$,
all from $x = \val{xOne}$ to the board's edge (no acceptance windows, kinetic windows or gates).
The player places up to four magnets of the listed strengths in the region
$\val{regX0} \le x \le \val{regX1}$, $\val{regY0} \le y \le \val{regY1}$. A flight ends when the
ion enters its detector, touches a magnet, or leaves the world (\xsecref{foundations}{sec:board}).

The reference places two magnets below the path, one of each orientation; together their fields
bend the ions upwards, the light ones most. That is the calutron's principle: in a given magnetic
field the curvature of a path is $qB/p$, and at equal energy the momentum grows as $\sqrt{m}$.
This document derives that relation exactly, as an integral along each path, and integrates the
three flights numerically.

\section{Setup}\label{sec:setup}

The level file, field by field (the appendix reproduces it):
\setuptable

In symbols: every ion has charge $q = \val{q}$ and radius $0$; $c = \val{c}$. With $T_0$ the same,
the momenta are $p = \val{pOne}$, $\val{pTwo}$ and $\val{pFour}$ (\xeqref{foundations}{eq:launch}),
nearly in the ratio $1 : \sqrt2 : 2$. The reference magnets are $\mu_1 = \val{muOne}$ at
$(\val{xMagOne}, \val{yMagOne})$ and $\mu_2 = \val{muTwo}$ at $(\val{xMagTwo}, \val{yMagTwo})$,
radius $\val{Rm}$ (field units, \physref{2.2}). The world is $-1 \le x \le \val{xWorld}$,
$-1 \le y \le \val{yWorld}$.

\begin{figure}[htbp]
\centering
\begin{tikzpicture}
\begin{axis}[axis equal image, width=0.92\linewidth, xmin=-1.6, xmax=31.6, ymin=-1.6, ymax=21.6,
  xtick={0,5,...,30}, ytick={0,5,...,20}, tick label style={font=\small},
  xlabel={$x$}, ylabel={$y$}, label style={font=\small}]
  \draw[black!8] (0,0) grid (30,20);
  \draw[black!45, dashed] (-1,-1) rectangle (\rawval{xWorld},\rawval{yWorld});
  \draw[black!30, dotted] (\rawval{regX0},\rawval{regY0}) rectangle (\rawval{regX1},\rawval{regY1});
  \fill[green!50!black, fill opacity=0.18] (\rawval{xOne},\rawval{yDetLoOne}) rectangle (30,\rawval{yDetHiOne});
  \draw[green!40!black] (\rawval{xOne},\rawval{yDetLoOne}) rectangle (30,\rawval{yDetHiOne});
  \fill[green!50!black, fill opacity=0.18] (\rawval{xOne},\rawval{yDetLoTwo}) rectangle (30,\rawval{yDetHiTwo});
  \draw[green!40!black] (\rawval{xOne},\rawval{yDetLoTwo}) rectangle (30,\rawval{yDetHiTwo});
  \fill[green!50!black, fill opacity=0.18] (\rawval{xOne},\rawval{yDetLoFour}) rectangle (30,\rawval{yDetHiFour});
  \draw[green!40!black] (\rawval{xOne},\rawval{yDetLoFour}) rectangle (30,\rawval{yDetHiFour});
  \addplot[thick, blue!70!black] table {pathone.dat};
  \addplot[thick, red!70!black] table {pathtwo.dat};
  \addplot[thick, orange!80!black] table {pathfour.dat};
  \fill[black!70] (\rawval{xMagOne},\rawval{yMagOne}) circle[radius=\rawval{Rm}];
  \fill[black!40] (\rawval{xMagTwo},\rawval{yMagTwo}) circle[radius=\rawval{Rm}];
  \fill (\rawval{xZero},\rawval{yZero}) circle[radius=0.15];
\end{axis}
\end{tikzpicture}
\caption{The two reference magnets below the path ($\otimes$ dark, $\odot$ light), the three
detectors (green), and the three flights: mass 1 (blue) bent most, into the top detector; mass 2
(red) into the middle one; mass 4 (orange) least, into the bottom one.}
\label{fig:board}
\end{figure}

\section{Principles and model}\label{sec:model}

\begin{principles}
  \item The magnetic force $q\vect{v}\times\vect{B}$ in the relativistic equation of motion
    \xeqref{foundations}{eq:eom}; it does no work (\xsecref{first_coil}{sec:nowork}).
  \item The magnets' fields in the plane, \xeqref{first_magnet}{eq:dipole-plane}, summed.
  \item The board's rules; the classical Runge--Kutta method, cubic Hermite interpolation, the
    error estimate by doubling the step and Simpson's rule (\xsecref{foundations}{sec:board},
    \xsecref{foundations}{sec:numerics}), as in \xsecref{slingshot}{sec:method}.
\end{principles}

The assumptions are First magnet's (\xsecref{first_magnet}{sec:model}): the magnets are held in
place, the ions are points without moments, and their radiation is not taken from them.

\section{Derivation}\label{sec:derivation}

\subsection{Curvature and rigidity}\label{sec:rigidity}

With only a magnetic field $B_z\uvec{z}$, $\abs{\vect{p}} = p$ is constant and
$\dot{\vect{p}} = qB_z(v_y, -v_x)$ turns $\vect{p}$ without changing its size. Write the direction
of motion as the angle $\theta$, $\vect{p} = p(\cos\theta, \sin\theta)$; then
$p\,\dot\theta = \hat{\vect{z}}\cdot(\vect{p}\times\dot{\vect{p}})/p = -qB_z v$, and with the arc
length $\dd s = v\,\dd t$,
\begin{equation}\label{eq:curvature}
  \deriv{\theta}{s} = -\frac{qB_z}{p}, \qquad
  \theta_\mathrm{e} = \theta_0 - \frac{q}{p}\int_\mathrm{path} B_z\,\dd s .
\end{equation}
The curvature is $qB_z/p$: the magnetic rigidity $p/q$ measures how hard an ion is to bend. In
the same field a lighter ion, of smaller $p$ at equal energy, turns more: by the factor
$\sqrt{m_2/m_1}$ if the paths sampled the same field. They do not, here: the paths part, and
each picks up its own $\int B_z\,\dd s$, so the angles differ by more than the momenta. The
relation \eqref{eq:curvature} itself is exact, and it is a check on each flight below.

\subsection{The method}\label{sec:method}

The three flights are integrated by \xeqref{foundations}{eq:rk4} at $h = \val{H}$ in the summed
field \xeqref{first_magnet}{eq:dipole-plane}, the events located on the cubic Hermite
interpolant. Along each, $\int B_z\,\dd s$ is summed by Simpson's rule on the steps to the exit
time, and \eqref{eq:curvature} gives the exit angle independently of the momentum the march
carries.

\section{The reference solution by hand}\label{sec:hand}

\begin{procedure}
  \item\label{st:launch} \emph{The launches.} The three momenta $\val{pOne}$, $\val{pTwo}$,
    $\val{pFour}$ along $+x$.
  \item\label{st:march} \emph{The march} and \emph{the bends.} Each flight, its
    $\int B_z\,\dd s$, the exit angle by \eqref{eq:curvature} and by the march, and the entry:
    \begin{center}\small\val{bends}\end{center}
  \item\label{st:entries} \emph{The entries.} Mass 1 enters its detector through the
    \val{faceOne} face at $(\val{xOne}, \val{yOne})$ at $t = \val{tOne}$; mass 2 at
    $(\val{xTwo}, \val{yTwo})$ at $t = \val{tTwo}$; mass 4 at $(\val{xFour}, \val{yFour})$ at
    $t = \val{tFour}$.
  \item\label{st:checks} \emph{The checks.} The two exit angles agree to $\val{bendAgree}$
    degrees; $\abs{\vect{p}}$ drifts by at most $\val{pDevOne}$; the step-doubling estimates are at
    most $\val{errMax}$.
\end{procedure}

\begin{result}[The reference flights]
The three isotopes enter their detectors at $(\val{xOne}, \val{yOne})$, $(\val{xTwo}, \val{yTwo})$
and $(\val{xFour}, \val{yFour})$. The level is solved.
\end{result}

\section{Comparison with the game}\label{sec:comparison}

\comparisontable

All differences are a few $10^{-11}$ or less, the game's verified accuracy (Foundations §6); our
estimates are below $10^{-12}$. The tolerance, $10^{-9}$, is ten times the larger of the two
computations' estimates, or the game's own preview--verification gap where that is larger
(mass 1), fixed before comparing.

\section{Other solutions}\label{sec:others}

On the empty board, the second placement in \code{placements.json}, the ions fly straight along
$y = \val{yZero}$: mass 1 leaves the world (\val{bareOutOne}), while masses 2 and 4 touch their
detectors' edges, which both lie on $y = \val{yZero}$ (\val{bareOutTwo}, \val{bareOutFour}). A
solution must lift mass 1 most and mass 4 least, which any field that curves the paths does by
\eqref{eq:curvature}; the art is in sending each into its own band. A map of solutions would take
the game's verdicts, which this document does not reproduce.

\section{What the game computes}\label{sec:engine}

The game takes each magnet's in-plane field as the dipole's (\physref{2.2}) and integrates the
equations of motion with its adaptive eighth-order Runge--Kutta method (\physref{5},
\physref{6}, \physref{7}); the differences in the comparison are the two integrations' errors.

\section{Further reading}\label{sec:reading}

Motion in a magnetic field and the radius $p/(qB)$: \Jackson{§12.2}. The magnets' field:
\levelref{first_magnet}. The calutron separated uranium isotopes at Oak Ridge by this principle,
with a nearly uniform field over a half-turn.

\levelfileappendix

\end{document}
